\section{Questão 2}
Encontre as PMFs marginais $P_X(x)$ e $P_Y(y)$ para o experimento randômico mostrado.

\begin{respostas}

As PMFs marginais são obtidas somando-se as linhas ($P_X$) e colunas ($P_Y$) da PMF conjunta. 
Em particular, como a PMF é simétrica, $P_X(X=a) = P_Y(Y=a)$.
Por inspeção da PMF no item (a),
%
\begin{align}
    P_X(X=0) &= \frac{4}{32} &
    P_X(X=1) &= \frac{10}{32} &
    P_X(X=2) &= \frac{14}{32} &
    P_X(X=3) &= \frac{4}{32} 
\\
    P_Y(Y=0) &= \frac{4}{32} &
    P_Y(Y=1) &= \frac{10}{32} &
    P_Y(Y=2) &= \frac{14}{32} &
    P_Y(Y=3) &= \frac{4}{32}
\end{align}

\end{respostas}